\documentclass[a4paper,12pt,twocolumn]{article}

% ---------- PACOTES ESSENCIAIS ----------
\usepackage[utf8]{inputenc}
\usepackage[T1]{fontenc}
\usepackage[brazil]{babel}
\usepackage{geometry}
\usepackage{amsmath, amssymb, amsfonts}
\usepackage{lmodern}
\usepackage{titlesec}
\usepackage{xcolor}
\usepackage[most]{tcolorbox}
\usepackage{graphicx}
\usepackage{fancyhdr}
\usepackage{hyperref}
\usepackage{setspace}

% ---------- CONFIGURAÇÃO DE PÁGINA E COLUNAS ----------
\geometry{top=2.5cm, bottom=2.5cm, left=1.8cm, right=1.8cm, columnsep=0.8cm}
\onehalfspacing

% ---------- CORES PERSONALIZADAS ----------
\definecolor{destaque}{RGB}{180,0,0} 

% ---------- CAIXA PARA DESTAQUES ----------
\newtcolorbox{caixadestaque}[1][]{
  enhanced,
  colback=destaque!5!white,
  colframe=destaque!80!black,
  boxrule=1pt,             
  left=0.5em, right=0.5em, top=0.5em, bottom=0.5em,
  sharp corners,
  title=#1,
  fonttitle=\bfseries
}

% ---------- CABEÇALHOS E RODAPÉS ----------
\pagestyle{fancy}
\fancyhf{}
\lhead{Projeto de Aplicações - 2026.1}
\rhead{Fundamentos da Álgebra Linear}
\rfoot{\thepage}

\fancypagestyle{firstpage}{
    \fancyhf{}
    \renewcommand{\headrulewidth}{0pt}
    \renewcommand{\footrulewidth}{0pt}
    \rfoot{\thepage}
}

% ---------- FORMATAÇÃO DOS TÍTULOS DE SEÇÃO ----------
\titleformat{\section}{\normalfont\large\bfseries\color{destaque}}{\thesection}{1em}{}
\titleformat{\subsection}{\normalfont\normalsize\bfseries\color{destaque!80!black}}{\thesubsection}{1em}{}

% =========================================================
% INFORMAÇÕES DO PROJETO (TÍTULO E AUTORES)
% =========================================================

\title{
    \vspace{-1.5cm}
    \includegraphics[width=0.4\textwidth]{ufpe_cin (1).png}\\[1em]
    \Large \textbf{Universidade Federal de Pernambuco -- UFPE}\\
    \large Centro de Informática -- CIn\\[0.5cm]
    \huge \textbf{Criptografia (Cifra de Hill)}
}

\author{
    \textbf{Gustavo Franco Pires} \\ \small gfp@cin.ufpe.br \and
    \textbf{João Azevedo Cavalcanti} \\ \small jac2@cin.ufpe.br \and
    \textbf{João Victor S. B. dos Santos} \\ \small jvsbs@cin.ufpe.br \and
    \textbf{Leonardo Portela Chiu} \\ \small lpc4@cin.ufpe.br
}
\date{\vspace{0.5em} \normalsize 2026} 

\begin{document}

\twocolumn[
  \begin{@twocolumnfalse}
    \maketitle
    
    \begin{abstract}
        O presente projeto aborda a aplicação de conceitos da Álgebra Linear na Criptografia, mais especificamente utilizando a Cifra de Hill. O método ilustra a utilização de matrizes, multiplicações e cálculo de determinantes para codificação e decodificação de mensagens através de transformações lineares no espaço modular. São detalhados os processos matemáticos de cifragem usando matrizes chaves invertíveis módulo 26, bem como a vulnerabilidade do método a ataques de texto plano conhecido, evidenciando o vínculo direto entre a teoria matemática estudada e sua aplicação prática no mundo real.
        
        \vspace{0.5em}
        \textbf{Palavras-chave:} Álgebra Linear, Criptografia, Cifra de Hill, Matrizes, Aritmética Modular.
    \end{abstract}
    \vspace{0.8cm}
  \end{@twocolumnfalse}
]

\thispagestyle{firstpage}

\section{Introdução à Criptografia de Hill}
\subsection{Introdução e Contexto Histórico}
A criptografia é a ciência e a arte de proteger a informação, transformando-a de modo a que apenas o destinatário legítimo, que possui o conhecimento de uma chave secreta, a consiga decifrar e ler. Desde a antiguidade, com métodos rudimentares como a Cifra de César, até aos complexos algoritmos digitais que protegem as transações bancárias atuais, a matemática tem sido o pilar fundamental da segurança da informação.

No campo da Álgebra Linear, destaca-se a Cifra de Hill, desenvolvida pelo matemático Lester S. Hill em 1929. Na época, a maioria das cifras conhecidas substituía uma letra de cada vez (cifras monoalfabéticas), o que as tornava vulneráveis à análise de frequência de letras num idioma. A Cifra de Hill foi pioneira ao introduzir uma cifra poligráfica prática, ou seja, um sistema que processa blocos de várias letras simultaneamente. Esta abordagem baseia-se inteiramente em operações com matrizes e aritmética modular, o que aumentou significativamente a segurança da época.

\subsection{O Alfabeto Numérico}
Para que seja possível aplicar ferramentas matemáticas (como matrizes e vetores) a um texto literário, o primeiro passo consiste em converter os caracteres alfabéticos em entidades numéricas. Utilizaremos o alfabeto padrão de 26 letras, mapeando cada caractere para um número inteiro de 0 a 25. Dessa forma, garantimos uma correspondência direta ideal para a aritmética modular de base 26:

\begin{center}
A = 0 | B = 1 | C = 2 | D = 3 | E = 4 | F = 5 | G = 6 \\
H = 7 | I = 8 | J = 9 | K = 10 | L = 11 | M = 12 | N = 13 \\
O = 14 | P = 15 | Q = 16 | R = 17 | S = 18 | T = 19 | U = 20 \\
V = 21 | W = 22 | X = 23 | Y = 24 | Z = 25
\end{center}

\subsection{Vetorização da Mensagem}
Após a conversão numérica, o texto original (também chamado de texto limpo) é segmentado em blocos de tamanho constante `n'. O valor de n define a dimensão do espaço vetorial com que vamos trabalhar e determina o tamanho da matriz quadrada que será usada como chave. Cada bloco de letras é estruturado na forma de um vetor-coluna $P \in \mathbb{R}^n$.

Exemplo Prático: Para ilustrar o funcionamento do nosso trabalho, escolhemos a palavra "HILL" e utilizaremos blocos de tamanho n = 2.

\begin{itemize}
    \item \textbf{Conversão de Letras para Números:} H $\rightarrow$ 7, I $\rightarrow$ 8, L $\rightarrow$ 11, L $\rightarrow$ 11
    \item \textbf{Agrupamento em Blocos e Criação dos Vetores-Coluna:} Como n = 2, dividimos as 4 letras em dois blocos emparelhados de duas letras:
    \begin{itemize}
        \item Bloco 1 ("HI"): Composto por H (7) e I (8).
        \item Bloco 2 ("LL"): Composto por L (11) e L (11).
    \end{itemize}
    \item \textbf{Montando os vetores-coluna matemáticos da nossa mensagem, obtemos:}
    \begin{equation}
        P_1 = \begin{pmatrix} 7 \\ 8 \end{pmatrix}, \quad P_2 = \begin{pmatrix} 11 \\ 11 \end{pmatrix}
    \end{equation}
\end{itemize}

\section{O Processo de Criptografia}
\subsection{Matriz Chave e seus Critérios}
Para que o processo de criptografia utilizando a Cifra de Hill funcione, o elemento central é a escolha de uma \textbf{matriz codificadora}, que chamaremos de Matriz Chave (A). Essa matriz representa a transformação matemática responsável por modificar a mensagem original e gerar o texto criptografado. A matriz A deve ser uma matriz quadrada de ordem NxN, como por exemplo uma matriz 2x2 ou 3x3, formada por números inteiros. O tamanho N dessa matriz dita o tamanho do vetor-coluna, ou seja, em blocos de quantas letras a mensagem original será dividida e cifrada por vez. Ou seja, uma matriz 2x2 irá criptografar a mensagem em grupos de duas letras por vez, enquanto uma matriz 3x3 irá trabalhar com grupos de três letras.
\begin{equation}
A = \begin{pmatrix} a_{11} & a_{12} \\ a_{21} & a_{22} \end{pmatrix}
\end{equation}

No entanto, a escolha da Matriz Chave não pode ser feita de forma aleatória. O critério mais rigoroso para a escolha da Matriz Chave é matemático: ela precisa obrigatoriamente ser \textbf{invertível módulo 26} (considerando o alfabeto padrão de 26 letras). Isso acontece, pois o processo de descriptografia depende da aplicação dessa matriz inversa para recuperar a mensagem original.

Para que a matriz possua uma inversa, o determinante da matriz A, denotado por det(A), não pode ser igual a zero. Entretanto, como a Cifra de Hill trabalha usando aritmética modular, essa condição precisa ser adaptada. Para que exista uma inversa módulo 26, é necessário que o máximo divisor comum (MDC) entre o valor de det(A) e o número 26 deve ser exatamente 1. Em termos práticos, isso significa que o determinante da Matriz Chave não pode ser um número par, múltiplo de 2, e nem múltiplo de 13. Essa condição garante que exista um inverso multiplicativo do determinante dentro da aritmética modular, tornando possível o processo inverso de decodificação.

\subsection{A Equação de Cifragem}
Uma vez estabelecida e validada a Matriz Chave, inicia-se a etapa prática de codificação. Primeiramente, o texto original, chamado de \textbf{texto plano}, é convertido para uma representação numérica utilizando uma tabela de correspondência entre letras e números (sendo a mais comum A=0, B=1, ..., Z=25). Depois da conversão, a mensagem é dividida em blocos de tamanho N, formando vetores-coluna chamados de \textbf{vetores de mensagem (P)}.

\begin{equation}
P = \begin{pmatrix} p_{1} \\ p_{2
\end{pmatrix}
\end{equation}

A cifragem ocorre estritamente através da operação multiplicação de matrizes. Para obter o vetor cifrado (C), multiplicamos a matriz chave A pelo vetor de mensagem P, baseando-se na seguinte equação algébrica:

\begin{equation}
C = A \cdot P
\end{equation}

Essa operação funciona como uma transformação linear: a matriz (A) altera matematicamente o vetor da mensagem original, gerando novos valores numéricos que escondem o conteúdo inicial. Na prática, cada elemento do vetor cifrado é obtido multiplicando os valores de uma linha da matriz pelos valores correspondentes da coluna do vetor mensagem e somando os resultados. O vetor obtido após essa multiplicação ainda não representa diretamente letras, pois alguns valores podem ultrapassar o intervalo permitido pelo alfabeto. Por isso, é aplicada a aritmética modular.

\subsection{O Papel da Aritmética Modular}
Ao realizarmos a multiplicação de matrizes $C = A \cdot P$, é esperado que os valores numéricos gerados no vetor C sejam altos, frequentemente ultrapassando a marca de 25, que representa a quantidade de letras disponíveis na nossa tabela alfabética. E para corrigir essa distorção e ciclar os números que introduzimos a aritmética modular. Aplicamos o conceito de \textit{módulo 26} aos elementos do vetor resultante C. Matematicamente, isso equivale a dividir cada número do vetor por 26 e extrair exclusivamente o resto dessa divisão inteira. Ao substituir os números brutos por seus respectivos restos, garantimos que todos os valores voltem a se enquadrar no intervalo de correspondência do nosso alfabeto. A partir desses novos números validados pelo módulo, as letras correspondentes são resgatadas, gerando o texto criptografado final.

\section{O Processo de Descriptografia}
\subsection{Inverso Multiplicativo}
Na criptografia, o seu objetivo ao descriptografar é isolar a mensagem original, eliminando o efeito da chave que a embaralhou. No momento que a letra original foi multiplicada, ou seja, precisamos achar essa relação:

Seja a um número inteiro pertencente ao conjunto de resíduos tal que $( a \cdot a^{-1}) = 1 \pmod m$. Temos que $a^{-1}$ é chamado de inversor multiplicativo ou recíproco.

Na aritmética modular, o recíproco deve ser obrigatoriamente um número inteiro que, multiplicado pelo número original, deixa o resto igual a 1 quando dividido pelo módulo m.

Ou seja, $( a \cdot a^{-1}) = (\text{múltiplo qualquer de } m) + 1$, implicando também que $(\text{múltiplo de } a) - (\text{múltiplo de } m) = 1$. Outra propriedade importante é que, assim como citado anteriormente, a ser invertível não pode ter primos em comum com o módulo m. Isso ocorre pois, se dois números compartilham uma base, qualquer combinação entre eles só consegue andar em "passos" do tamanho dessa base, nunca de 1 em 1.

Exemplo: Seja $m = 26$, então seus fatores primos são 2 e 13.
$a = 4$ (múltiplo de 2)
(múltiplo de 4) - (múltiplo de 26) = $2 \cdot [(\text{múltiplo } 2) - (\text{múltiplo de } 13)]$. O resultado nunca resultará em 1. Isso vale assim como para os múltiplos de 13.

\subsection{Matriz inversa}
Podemos transformar $C = A \cdot P$ em $P = A^{-1} \cdot C \pmod{26}$.

Para revelar a mensagem original, o processo é invertido: o vetor criptografado deve ser multiplicado pela matriz inversa módulo 26 ($A^{-1}$). Na Criptografia de Hill nós não podemos usar divisões nem frações, porque precisamos de números inteiros para representar letras. A base do processo usado é achar a inversa de ($A^{-1}$) por matriz adjunta e adaptá-la para a aritmética modular.

A teoria subjacente à inversão matricial fundamenta-se na busca pela Matriz Identidade, o elemento neutro da álgebra linear, definida pela relação $A \cdot A^{-1} = I$. Pelo Teorema de Laplace, sabe-se que o produto de uma matriz pela sua respectiva Matriz Adjunta resulta em uma matriz escalar cujos elementos da diagonal principal são iguais ao determinante da matriz original ($A \cdot \text{Adj}(A) = \det(A) \cdot I$). Para isolar a Identidade, a álgebra clássica exige a divisão escalar da equação pelo determinante. No entanto, na modelagem da Cifra de Hill, devido ao domínio restrito do alfabeto criptográfico, a operação de divisão é substituída pela multiplicação do inverso multiplicativo modular do determinante. Essa adaptação teórica garante a reversibilidade da matriz adjunta sem a introdução de valores fracionários, mantendo todo o sistema contido no anel de inteiros módulo 26.

Inicialmente a fórmula para achar a matriz inversa por adjunta é essa:
\begin{equation}
A^{-1} = \frac{1}{\det(A)} \cdot \text{Adj}(A)
\end{equation}

Porém, como em aritmética modular não permite frações, iremos substituir pelo recíproco modular do determinante, denotado por $\det(A)^{-1}$. Logo depois, multiplicarmos pela $\text{Adj}(A)$ com o resíduo do mod(26):
\begin{equation}
A^{-1} = \det(A)^{-1} \cdot \text{Adj}(A) \pmod{26}
\end{equation}

Se o resíduo de $\det(A) = ab - cd$ módulo de 26 não for divisível por 13 ou 2, então a sua inversa é dada por:
\begin{equation}
A^{-1} = (ad-bc)^{-1} \begin{pmatrix} d & -b \\ -c & a \end{pmatrix} \pmod{26}
\end{equation}

Uma excelente prática é sempre que achar a matriz inversa, verificar o seguinte teorema:
\begin{equation}
A \cdot A^{-1} = I \pmod{26}
\end{equation}

Após achar a matriz inversa o processo é resumidamente igual ao processo de criptografar:
\begin{equation}
P = A^{-1} \cdot C \pmod{26}
\end{equation}

\section{Exemplo Prático Completo}
Com base em tudo que aprendemos até agora sobre Criptografia, sobretudo, Cifra de Hill, vamos exemplificar para ver se entendemos de fato.

Exemplo: O professor Silvio queria testar se os seus alunos estavam dominando Cifra de Hill e ele passou 2 atividades para eles.

Observações:
\begin{itemize}
    \item O alfabeto baseado vai seguir o padrão: A=1, B=2, \dots, Y=25, Z=0
    \item As atividades abaixo deverão utilizar como matriz base a matriz abaixo:
\end{itemize}
\begin{equation}
A = \begin{pmatrix} 5 & 6 \\ 2 & 3 \end{pmatrix}
\end{equation}

a) Cifrar a palavra MATEMATICA, tendo como base o que foi aprendido

b) Decifrar a mensagem YLMCNFSV, tendo como base o que foi aprendido

\subsection{Resolução das Atividades}
\textbf{Atividade a)}

Primeiro, antes de tudo, vamos converter as letras da palavra MATEMATICA em números.
M = 13; A = 1; T = 20; E = 5; M = 13; A = 1; T = 20; I = 9; C = 3; A = 1 (Nota: a resolução assume a conversão das letras finais)

Como se trata de uma matriz 2x2, então será uma Cifra de Hill de ordem 2. Logo, agrupamos as letras em pares. Cada par formará um vetor coluna como abaixo:
\begin{equation}
\begin{pmatrix} 13 \\ 1 \end{pmatrix}, \begin{pmatrix} 20 \\ 5 \end{pmatrix}, \begin{pmatrix} 13 \\ 1 \end{pmatrix}, \begin{pmatrix} 20 \\ 9 \end{pmatrix}, \begin{pmatrix} 3 \\ 1 \end{pmatrix}
\end{equation}

Após formar os vetores coluna, o próximo passo da resolução é multiplicar a matriz chave pelos vetores coluna para achar os vetores coluna cifrados. Após as multiplicações de matrizes, aplique o módulo 26.

\begin{equation}
\begin{pmatrix} 5 & 6 \\ 2 & 3 \end{pmatrix} \cdot \begin{pmatrix} 13 \\ 1 \end{pmatrix} = \begin{pmatrix} 71 \\ 29 \end{pmatrix} \rightarrow \begin{pmatrix} 19 \\ 3 \end{pmatrix} \rightarrow \text{SC}
\end{equation}

\begin{equation}
\begin{pmatrix} 5 & 6 \\ 2 & 3 \end{pmatrix} \cdot \begin{pmatrix} 20 \\ 5 \end{pmatrix} = \begin{pmatrix} 130 \\ 55 \end{pmatrix} \rightarrow \begin{pmatrix} 0 \\ 3 \end{pmatrix} \rightarrow \text{ZC}
\end{equation}

\begin{equation}
\begin{pmatrix} 5 & 6 \\ 2 & 3 \end{pmatrix} \cdot \begin{pmatrix} 13 \\ 1 \end{pmatrix} = \begin{pmatrix} 71 \\ 29 \end{pmatrix} \rightarrow \begin{pmatrix} 19 \\ 3 \end{pmatrix} \rightarrow \text{SC}
\end{equation}

\begin{equation}
\begin{pmatrix} 5 & 6 \\ 2 & 3 \end{pmatrix} \cdot \begin{pmatrix} 20 \\ 9 \end{pmatrix} = \begin{pmatrix} 154 \\ 67 \end{pmatrix} \rightarrow \begin{pmatrix} 24 \\ 15 \end{pmatrix} \rightarrow \text{XO}
\end{equation}

\begin{equation}
\begin{pmatrix} 5 & 6 \\ 2 & 3 \end{pmatrix} \cdot \begin{pmatrix} 3 \\ 1 \end{pmatrix} = \begin{pmatrix} 21 \\ 9 \end{pmatrix} \rightarrow \begin{pmatrix} 21 \\ 9 \end{pmatrix} \rightarrow \text{UI}
\end{equation}

Portanto, a resposta cifrada da atividade A ficará: SCZCSCXOUI

\vspace{0.5cm}
\textbf{Atividade b)}

Para decifrar a mensagem YLMCNFSV, temos que fazer o processo inverso: encontrar a matriz inversa em módulo 26 e multiplicá-la pelos vetores cifrados.

Calculando o determinante de A:
det(A) = (5 . 3) - (6 . 2) = 15 - 12 = 3

Para achar a matriz inversa, precisamos multiplicar a matriz adjunta pelo inverso do determinante. Porém, como não são aceitas frações em aritmética modular, ao invés disso, iremos substituir pelo recíproco modular do det(A). No caso dessa atividade, o recíproco modular é 9.

Agora, vamos montar a matriz adjunta assim:
\begin{equation}
\text{Adj}(A) = \begin{pmatrix} 3 & -6 \\ -2 & 5 \end{pmatrix} \pmod{26} = \begin{pmatrix} 3 & 20 \\ 24 & 5 \end{pmatrix}
\end{equation}

Então, a matriz inversa será o resultado da multiplicação entre o recíproco modular do det(A) e a adj(A):
\begin{equation}
A^{-1} = 9 \cdot \begin{pmatrix} 3 & 20 \\ 24 & 5 \end{pmatrix} = \begin{pmatrix} 27 & 180 \\ 216 & 45 \end{pmatrix} = \begin{pmatrix} 1 & 24 \\ 8 & 19 \end{pmatrix}
\end{equation}

Logo após achar a matriz inversa, o próximo passo é converter as letras da mensagem criptografada em números e depois separá-las em pares de vetores:
Y = 25; L = 12; M = 13; C = 3; N = 14; F = 6; S = 19; V = 22;

Após formar os vetores coluna, o próximo passo da resolução é multiplicar a matriz inversa pelos vetores cifrados para achar os vetores originais. Após as multiplicações de matrizes, aplique o módulo 26.

\begin{equation}
\begin{pmatrix} 1 & 24 \\ 8 & 19 \end{pmatrix} \cdot \begin{pmatrix} 25 \\ 12 \end{pmatrix} = \begin{pmatrix} 313 \\ 428 \end{pmatrix} \rightarrow \begin{pmatrix} 1 \\ 12 \end{pmatrix} \rightarrow \text{AL}
\end{equation}

\begin{equation}
\begin{pmatrix} 1 & 24 \\ 8 & 19 \end{pmatrix} \cdot \begin{pmatrix} 13 \\ 3 \end{pmatrix} = \begin{pmatrix} 85 \\ 161 \end{pmatrix} \rightarrow \begin{pmatrix} 7 \\ 5 \end{pmatrix} \rightarrow \text{GE}
\end{equation}

\begin{equation}
\begin{pmatrix} 1 & 24 \\ 8 & 19 \end{pmatrix} \cdot \begin{pmatrix} 14 \\ 6 \end{pmatrix} = \begin{pmatrix} 158 \\ 226 \end{pmatrix} \rightarrow \begin{pmatrix} 2 \\ 18 \end{pmatrix} \rightarrow \text{BR}
\end{equation}

\begin{equation}
\begin{pmatrix} 1 & 24 \\ 8 & 19 \end{pmatrix} \cdot \begin{pmatrix} 19 \\ 22 \end{pmatrix} = \begin{pmatrix} 547 \\ 570 \end{pmatrix} \rightarrow \begin{pmatrix} 1 \\ 24 \end{pmatrix} \rightarrow \text{AX}
\end{equation}

Portanto, a mensagem decifrada é ALGEBRAX (o X aparece pois, quando o número de letras em uma Cifra de Hill de ordem 2 é ímpar, temos que colocar uma letra pra completar os pares formados nos vetores coluna).

\section{Conclusão}
\textbf{Vantagens:}
A Criptografia (Cifra de Hill) destaca-se pela sua capacidade de mascarar os padrões de linguagem. Isso torna o método altamente resistente aos ataques clássicos de análise de frequência de letras, que quebravam facilmente cifras mais antigas. Além disso, aumenta de forma significativa a complexidade do sistema.

\textbf{Desvantagens:}
Por outro lado, o método apresenta algumas desvantagens e vulnerabilidades matemáticas. A primeira limitação diz respeito à criação da chave: a matriz escolhida precisa ser obrigatoriamente invertível no módulo 26 (ou seja, seu determinante não deve ter fatores primos em comum com 26), o que restringe as opções de chaves válidas e exige cuidado na sua formulação. A falha mais crítica, no entanto, é a sua total vulnerabilidade ao Known-Plaintext Attack (KPA). Por ser um sistema puramente linear, um invasor que descubra um pequeno trecho da mensagem original e o seu correspondente cifrado consegue, através de cálculos básicos de matriz inversa, descobrir a chave secreta e comprometer todo o sistema, dependendo do tamanho da matriz chave.

Em um cenário de análise criptográfica, um atacante pode obter uma parte da mensagem original juntamente com a sua respectiva versão criptografada. Isso pode acontecer quando existem trechos previsíveis de uma comunicação ou quando parte da mensagem é conhecida. O objetivo do ataque é descobrir a matriz A, que é justamente a transformação linear responsável por realizar a criptografia.

Considerando uma Matriz Chave de ordem 2x2, temos:
\begin{equation}
A = \begin{pmatrix} a_{11} & a_{12} \\ a_{21} & a_{22} \end{pmatrix}
\end{equation}

Essa matriz possui quatro valores desconhecidos. Portanto, é necessário obter informações suficientes da mensagem original e da mensagem cifrada para determinar esses valores. Para isso, o atacante reúne blocos conhecidos da mensagem original e organiza esses valores em uma matriz P. Da mesma forma, os blocos correspondentes da mensagem criptografada são organizados em uma matriz C. Partindo da equação original, isolando a Matriz Chave utilizando a inversa da matriz P, chegamos a seguinte equação:
\begin{equation}
A = C \cdot P^{-1} \pmod{26}
\end{equation}

Ou seja, conhecendo um conjunto suficiente de pares entre texto original e texto criptografado, é possível recuperar a Matriz Chave através de operações de Álgebra Linear. Para que esse cálculo seja possível, a matriz formada pelo texto original precisa possuir uma inversa no módulo 26. Isso significa que o determinante da matriz P não pode ser múltiplo de 2 nem de 13. Caso essa condição não seja satisfeita, aquele conjunto específico de letras não permite encontrar a matriz chave. Porém, outro agrupamento de letras pode gerar uma matriz P diferente e possibilitar a recuperação da chave.

Em resumo, o estudo da Criptografia (Cifra de Hill) ilustra de forma clara os conceitos fundamentais de Álgebra Linear, como operações com matrizes, determinantes e inversão podem ser aplicados diretamente na área da criptografia, através do uso da aritmética modular.

\begin{thebibliography}{9}

\bibitem{anton} 
ANTON, Howard; RORRES, Chris. 
\textit{Álgebra Linear com Aplicações}. 
10ª Edição, Bookman, 2012.

\bibitem{strang} 
STRANG, Gilbert. 
\textit{Introdução à Álgebra Linear}. 
4ª Edição, Wellesley-Cambridge Press, 2016.

\end{thebibliography}

\end{document}
